Throughout this section $\mathcal U=\mathbb R^p$ denotes the (discretized) input
space and $\mathcal V=\mathbb R^q$ the output space, $G:\mathcal U\to\mathcal V$
the target operator, $\mu$ the input distribution, and
$(u_1,v_1),\dots,(u_N,v_N)$ i.i.d.\ draws with $v_n=G(u_n)$; the data carry no
sampling noise in the usual sense, coming from a deterministic solver, though
Section~\ref{sec:diversity} discusses solver error. Assume $\|G(u)\|_2>0$ almost surely and that all displayed expectations are finite. We write
$\ell(w,v)=\|w-v\|_2/\|v\|_2$ for the relative error and
$R(\widehat G)=\mathbb E_{u\sim\mu}\,\ell(\widehat G(u),G(u))$ for the risk, which is
the quantity reported in Section~\ref{sec:experiments}. The results below are
largely elementary, but each one motivates a specific design decision in the
pipeline of Section~\ref{sec:method}, and each has a measurable footprint in
the experiments: one statement per stage, from the symmetry handling and the
kernel-flow term through stacking, the residual correction, its coverage
guarantee, and the choice of nugget. Proofs are collected in
Supplement~\ref{app:proofs}.

\subsection{Guarantees for the fitted stages}
\label{sec:theory-stages-summary}

Every fitted stage of the pipeline has an elementary finite-sample
bound, stated and proved in the online supplement (Section~S1). Their
status should be read exactly: they are conditional capacity calculations
for candidates fixed before the validation sample is drawn, and the
pipeline as deployed does not satisfy that hypothesis, because the same
validation split that scores the stacks and the kernel grids also selects
the network checkpoints, the member set and, on OCO-2, the coordinate
winners. We therefore present them as the scaling behind each design
decision, with constants measured on the runs, and not as guarantees for
the shipped estimator; as Section~S1 notes, several are numerically vacuous
at the campaign's sizes in any case. Split-conformal coverage has a different independence requirement: the predictor and scale must be fixed independently of the calibration and evaluation data. The retrospective experiment is interpreted under that requirement, not as protection against unrestricted reuse of its public test block. The bounds are: a symmetrization inequality behind the reflection
averaging (Proposition~\ref{prop:sym}), the leave-half-out reading of
the kernel-flow regularizer (Proposition~\ref{prop:kf}), capacity
bounds for the validated stack, the honesty comparison, the
per-coordinate affine stack, and the grid-tuned correction
(Propositions~\ref{prop:stack} and~\ref{prop:affine},
Lemma~\ref{lem:select}, Corollary~\ref{cor:certificate}), the last
with an analytic uniform bound on the correction weights
(Lemma~\ref{lem:kappa}, $\kappa_{\mathrm{an}}=500$, sharp over kernels and
designs). The statistical capacity bounds are reported for their scaling. Theorems~\ref{thm:or} and~\ref{thm:minimax} give deterministic conditional bounds, whose unknown norm factors prevent numerical error certification.
\subsection{What the stack can achieve}
\label{sec:theory-secmom}

Proposition~\ref{prop:stack} bounds weight selection under independent-validation assumptions; it does not imply that mixing improves on a member. That is decided by a single measurable matrix. For a
map $f$ write the \emph{normalized residual} at $u$ as
$\rho_f(u)=\big(f(u)-G(u)\big)/\|G(u)\|_2$, so that
$\ell(f(u),G(u))=\|\rho_f(u)\|_2$.

\begin{proposition}[second-moment identity for convex stacks]\label{prop:secmom}
Let $f_1,\dots,f_M$ be fixed maps and define
$S\in\mathbb R^{M\times M}$ by
$S_{mk}=\mathbb E_{u\sim\mu}\,\langle\rho_{f_m}(u),\rho_{f_k}(u)\rangle$.
For $w$ in the simplex $\Delta^{M-1}$ let $f_w=\sum_m w_mf_m$. Then
\[
\mathbb E_{u\sim\mu}\,\ell\big(f_w(u),G(u)\big)^2\;=\;w^\top S\,w ,
\]
so the squared-metric risk of every convex stack is determined by $S$, and
the best achievable is $\min_{w\in\Delta^{M-1}}w^\top Sw$. In particular:
\begin{enumerate}
\item[(i)] (two members) if $S$ has diagonal $(e_1^2,e_2^2)$ with $0<e_1\le e_2$
and correlation $\varrho=S_{12}/(e_1e_2)$, mixing in the weaker member
strictly helps if and only if $\varrho<e_1/e_2$, in which case the optimum is
interior with value $e_1^2e_2^2(1-\varrho^2)/(e_1^2+e_2^2-2\varrho e_1e_2)$;
\item[(ii)] (equicorrelated family) if $S_{mm}=e^2$ and $S_{mk}=\varrho\,e^2$
for all $m\neq k$ with $\varrho\ge0$, then
$\min_{w}w^\top Sw=e^2\big(\varrho+(1-\varrho)/M\big)$, attained at uniform
weights and decreasing to $e^2\varrho$ as $M\to\infty$.
\end{enumerate}
\end{proposition}
\noindent\emph{Proof.} See Supplement~\ref{proof:prop:secmom}.


If a member has zero squared risk, it already attains the minimum zero; ratios involving its error are not used.

Part~(i) is the classical two-member minimum-variance rule, and part~(ii) is
the equicorrelated case of the ambiguity decomposition of ensemble learning
\citep{krogh1995neural}; see \citet{ueda1996generalization} for the
generalization-error form and \citet{brown2005diversity} for a survey. Both
enter here as measurable decision tools on the estimated $S$.

Everything in this subsection is stated for the squared metric, and the tables
of Sections~\ref{sec:experiments} and~\ref{sec:oco2} report a different one.
The distinction is small numerically and decisive logically, so we make it
explicit once and refer back to it wherever the two families of numbers meet.

\begin{remark}[the two error metrics]\label{rem:metric}
For a map $f$ write
$\mathcal E_1(f)=\mathbb E\|\rho_f(u)\|_2$ and
$\mathcal E_2(f)=(\mathbb E\|\rho_f(u)\|_2^2)^{1/2}$, the mean and the root
mean square of the per-sample relative error. If $\mathcal E_1=0$, both errors vanish and no ratio is needed; below assume $\mathcal E_1>0$. The tables report
$\mathcal E_1$; Proposition~\ref{prop:secmom} is a statement about
$\mathcal E_2$, since $\mathcal E_2(f_w)^2=w^\top Sw$ by definition. Writing
$X=\|\rho_f(u)\|_2$ and $c_f=(\operatorname{Var}X)^{1/2}/\mathbb EX$,
\[
\mathcal E_2(f)^2-\mathcal E_1(f)^2=\operatorname{Var}X,
\qquad
\mathcal E_1(f)/\mathcal E_2(f)=(1+c_f^2)^{-1/2}\le1 ,
\]
with equality only if the per-sample error is almost surely constant. Two
consequences are used in this paper. The floor belongs in $\mathcal E_2$,
because only $\mathcal E_2$ is determined by $S$: the map
$w\mapsto w^\top Sw$ is a quadratic form in second moments, whereas
$w\mapsto\mathbb E\|\sum_m w_m\rho_{f_m}(u)\|_2$ depends on the full joint law
of the residuals. And upper bounds transfer downward for free, since
$\mathcal E_1\le\mathcal E_2$, while lower bounds do not --- so a floor read
off $S$ says nothing about a table entry without the dispersion of the same
predictor. Comparing the two directly is biased in favour of the table entry
by $(1+c_f^2)^{-1/2}$, which is about six percent on the structural-mechanics
stack and ten to eleven percent on the O2 objects of
Section~\ref{sec:oco2}; the dispersion is a property of the predictor and not
a constant of the problem, so a single conversion factor should not be reused
across objects.
\end{remark}

Part~(i) is an exact criterion, but it is used as a decision rule and a
decision rule is applied to estimates, so what governs it in practice is
not the criterion but the gap it has to resolve relative to the noise in
that estimation. A margin law for threshold rules of this shape
(Proposition~\ref{prop:margin}, Supplement~\ref{supp:rules}) bounds the
probability that a decision is both wrong and taken at an observed margin
of at least $\tau$ by $\exp(-\tau^{2}/2\sigma^{2})$, with $\sigma$ the
scale of the estimation error and not a chosen constant. That is a joint
bound, not the error rate among the decisions whose observed margin
exceeds $\tau$: the conditional quantity is the joint one divided by
$\mathbb P(|\widehat D|\ge\tau)$, on which the proposition says nothing.
Accuracy inside a margin bin is therefore an empirical calibration of the
rule, and Section~\ref{sec:secmom-margin} reports it that way.

Exact equicorrelation never holds in practice, but for convex weights the
floor survives one-sided bounds on the entries of $S$, which are what one
actually measures.

\begin{corollary}[ensembling floor]\label{cor:floor}
If every member satisfies $S_{mm}\ge\bar e^{\,2}$ and every pair satisfies
$S_{mk}\ge\bar\varrho\,\bar e^{\,2}$ with $\bar\varrho\in[0,1]$, then every
convex combination $f_w$ obeys
\[
\mathbb E\,\ell\big(f_w(u),G(u)\big)^2\;=\;w^\top Sw\;\ge\;
\bar e^{\,2}\big(\bar\varrho+(1-\bar\varrho)\,\|w\|_2^2\big)\;\ge\;
\bar e^{\,2}\,\bar\varrho .
\]
No number of additional members with the same error level and mutual
correlations moves the ensemble's root-mean-square relative error below
$\bar e\sqrt{\bar\varrho}$.
\end{corollary}
\noindent\emph{Proof.} See Supplement~\ref{proof:cor:floor}.


The floor is a lower bound; a matching upper bound follows from the same
entrywise information, and together they pin the achievable window from both
sides by measured quantities alone.

\begin{corollary}[the floor is nearly achieved]\label{cor:floorpinch}
Suppose $\bar e>0$ and, in addition to the hypotheses of Corollary~\ref{cor:floor}, that
$S_{mm}\le\bar e^{\,2}(1+\delta_e)$ for every $m$ and
$S_{mk}\le(\bar\varrho+\delta_\varrho)\,\bar e^{\,2}$ for every $m\neq k$.
Then
\[
\bar\varrho\;\le\;\frac{1}{\bar e^{\,2}}\min_{w\in\Delta^{M-1}}w^\top Sw
\;\le\;\bar\varrho+\delta_\varrho+\frac{1+\delta_e-\bar\varrho-\delta_\varrho}{M}.
\]
The width of the bracket is the measured correlation spread
$\delta_\varrho$ plus a $1/M$ term: when these entrywise spreads are small, the optimal convex risk is localized. The upper bound is obtained by evaluating uniform weights; it does not upper-bound an arbitrary fitted convex combination.
\end{corollary}
\noindent\emph{Proof.} See Supplement~\ref{proof:cor:floorpinch}.


Applied to a measured matrix $\widehat S$, these are exact empirical bounds on that same design, not confidence bounds on the population matrix. For the six-member validation matrices, the weak RMS floor $\sqrt{\min_{i\ne j}\widehat S_{ij}}$ is $4.820\pm0.062\%$, the stronger finite-six-member lower bound is $4.849\pm0.060\%$, and the simplex optimum is $4.920\pm0.057\%$ (mean and sample standard deviation over ten seeds). The values near $4.91\%$ obtained from mean centered correlations are equicorrelation diagnostics, not these lower bounds.

The floor is stated for convex weights, and the deployed stack is not
convex: its per-pixel weights are affine and unconstrained in sign. For
global weights the sign constraint can be removed exactly.

\begin{corollary}[signed stacks]\label{cor:signed}
If $S$ is positive definite, the minimum of $w^\top Sw$ over all $w$ with
$\sum_m w_m=1$, of either sign, is $1/(\mathbf 1^\top S^{-1}\mathbf 1)$,
attained at $w^\star=S^{-1}\mathbf 1/(\mathbf 1^\top S^{-1}\mathbf 1)$; it
is at most the convex minimum, with equality exactly when $w^\star$ has no
negative coordinate.
\end{corollary}
\noindent\emph{Proof.} See Supplement~\ref{proof:cor:signed}.


On the measured $S$ the two minima agree to within $0.0005$ points at every
seed and for both member sets, but not because the constraint is slack: the
unconstrained minimizer carries a small negative coordinate at seven seeds of
ten for the six members and at one seed for the five (the stored
\texttt{signed} records), and removing it costs at most those five
ten-thousandths. Allowing signed global weights therefore buys nothing
measurable here, and what the deployed per-pixel stack gains over the convex
one ($0.032$ and $0.035$ points in the two campaigns) is a benefit of a different per-pixel affine class; the global signed-weight comparison does not isolate the effect of signs or intercepts within that class.

The members of a pool are usually of two kinds at once, architectures and
seeds of each architecture, and the floor separates the two.

\begin{theorem}[seeds and architectures]\label{thm:blockfloor}
Let the pool consist of $K$ seeds of each of $M$ architectures, and suppose
$S_{ii}\ge\bar e^{\,2}$ for every member, $S_{ij}\ge\varrho_w\bar e^{\,2}$
for two seeds of the same architecture and $S_{ij}\ge\varrho_b\bar e^{\,2}$
for members of different architectures, with $0\le\varrho_b\le\varrho_w\le1$.
Then every convex combination obeys
\[
\mathbb E\,\ell\big(f_w(u),G(u)\big)^2\;\ge\;\bar e^{\,2}\Big(\varrho_b
+\frac{\varrho_w-\varrho_b}{M}+\frac{1-\varrho_w}{MK}\Big),
\]
with equality at uniform weights when the three inequalities are equalities.
For an extended pool only if the same entrywise bounds continue to hold, no number of seeds of one architecture takes the
root-mean-square error below $\bar e\sqrt{\varrho_w}$, no number of seeds of
$M$ architectures takes it below
$\bar e\sqrt{\varrho_b+(\varrho_w-\varrho_b)/M}$, and no pool of any size
takes it below $\bar e\sqrt{\varrho_b}$.
\end{theorem}
\noindent\emph{Proof.} See Supplement~\ref{proof:thm:blockfloor}.


At fixed $M$, increasing $K$ reduces only the term $(1-\varrho_w)/(MK)$ in the lower bound. The within-architecture contribution $(\varrho_w-\varrho_b)/M$ persists. Section~\ref{sec:seeds-arch} applies the formula empirically to sixty trained predictors; it does not assume the same inequalities for untrained members.

The two-member rule of Proposition~\ref{prop:secmom}(i) can be read the same
way, as a statement about what a member's accuracy is worth against its
correlation, and the exchange rate is explicit.

\begin{proposition}[the price of accuracy in correlation]\label{prop:exchange}
In the setting of Proposition~\ref{prop:secmom}(i) write $t=e_2/e_1$ and
assume $-1<\varrho<1/t$ (with $t\ge1$), so the optimum is interior. The optimal
two-member risk is $V=e_1^2\,t^2(1-\varrho^2)/(1+t^2-2\varrho t)$, and
\[
\frac{\partial V}{\partial t}
=\frac{2e_1^2\,t\,(1-\varrho^2)(1-\varrho t)}{(1+t^2-2\varrho t)^2}>0,
\qquad
\frac{\partial V}{\partial\varrho}
=\frac{2e_1^2\,t^2(t-\varrho)(1-\varrho t)}{(1+t^2-2\varrho t)^2}>0 .
\]
A change $(\mathrm dt,\mathrm d\varrho)$ in the second member therefore
leaves the optimal risk unchanged to first order exactly when
$\mathrm d\varrho=-(1-\varrho^2)\,\mathrm dt/\big(t(t-\varrho)\big)$: an
improvement of the second member's error by a fraction $\epsilon$ is
cancelled by a rise of $\epsilon(1-\varrho^2)/(t-\varrho)$ in its correlation
with the first.
\end{proposition}
\noindent\emph{Proof.} See Supplement~\ref{proof:prop:exchange}.


At perfect anticorrelation $\varrho=-1$, cancellation gives zero optimal risk and the derivative in $t$ is zero, so that boundary is excluded above. At the admission threshold $\varrho=1/t$ with $t>1$, the exchange coefficient $(1-\varrho^2)/(t-\varrho)$ tends to $1/t$; it need not be large. The formula is a local sensitivity calculation at an interior point. Supplement~\ref{supp:experiments} evaluates it at the measured values.

For coordinate-dependent combinations, squared losses with a fixed diagonal metric separate by coordinate (Proposition~\ref{prop:percoord}; proof in Supplement~\ref{app:proofs}). The same holds after a fixed sample weighting (Corollary~\ref{cor:percoordw}; proof in Supplement~\ref{app:proofs}). These statements motivate the selector but do not make it simultaneously optimal for the two reported mean relative norms: the square root couples coordinates, and the two criteria also have different sample denominators. Proposition~\ref{prop:selectcoord} gives an independent-validation selection bound (proof in Supplement~\ref{app:proofs}); the deployed candidates reuse validation, as noted above. The OCO-2 combination improves on the rescored source emulator in both mean criteria, while the flat feature head is slightly better on the O2 reduced criterion.

\subsection{Metric and representation: when a fixed kernel loses to a network}
\label{sec:theory-aniso}

The structural-mechanics benchmark has the neural mean and the isotropic
kernel within half a point of each other. On the OCO-2 emulation problem of
Section~\ref{sec:oco2} the same fixed kernel is an order of magnitude behind
the same network. Two mechanisms can produce such a gap, and they separate
cleanly.

The first is input geometry. A sensitivity-based linear transformation improves the raw-input kernel by a factor of about $1.1$--$1.5$ in this study. The learning-curve sweep does not establish a general rate improvement. Proposition~\ref{prop:aniso} in Supplement~\ref{supp:rules} gives a finite-design bound for a Lipschitz target after a rank-$r$ truncation, with its approximation error explicit; its proof is in Supplement~\ref{app:proofs}. It makes no assertion that a measured approximate rank changes an asymptotic rate.

The second is the representation used by the kernel. Native-space bounds depend on the target's norm in that space. A fixed kernel may approximate targets outside its RKHS, so native-space membership is not a limit on approximation expressivity. Learned features can change the relevant norm and design geometry, but neither change is automatically favorable.

The representational term has a precise reading through the same
optimal-recovery bound that Section~\ref{sec:theory-or} makes exact. That
bound controls the error by $\|G-m\|_{\mathcal H}\cdot P_\lambda$, a product of
the target's norm in the kernel's native space and a design factor determined by the kernel, training inputs, and query. Changing the features changes both in principle, and on the
O2 band the two move by very different amounts. The effective dimension of
the Mat\'ern Gram, the quantity Lemma~\ref{lem:deff} controls, is nearly
identical on the raw input and on the learned features (at $n=4000$ and a
matched median length scale, $\alignDeffRaw$ against $\alignDeffFeat$ at a
$10^{-8}$ nugget); the design factor itself is not fixed by that trace, and
measured directly at the $2000$ test inputs the feature kernel's power
function is smaller at the median scale, its median $\alignRatioP$ times
below the raw kernel's, and equal to it ($\alignRatioPSel$) at the
multiplier $4$ that the campaign selected for every feature head. The
fitted-norm diagnostic moves far more: regressing the same outputs through
the two kernels, the squared native-space norm of the fitted interpolant,
$\operatorname{tr}(\alpha^\top K\alpha)$, is $\alignRatioNorm$ times
smaller on the features than on the raw input at the median scale
($\alignRatioFull$ times on the full $18000$-row block) and
$\alignRatioSel$ times smaller at the selected multiplier ($\alignRatioFullSel$ times on the
full block). (An earlier
draft quoted $42$ from $\alpha^\top Y$, which adds the term
$n\lambda\|\alpha\|^2$; at this nugget the two agree to three digits, and
the network retrained for the check lands at $38$.) Both numbers are
finite-design diagnostics of the interpolating regime: for a target in the corresponding RKHS with compatible labels, the fitted norm is a lower bound on its norm, the ratio grows with the kernel
scale --- $\alignRatioLo$ at half the median length scale, $\alignRatioHi$
at twice it --- and at a $10^{-4}$ nugget, where the ridge no longer
interpolates, it collapses and at the selected multiplier inverts
($\alignRatioHeavy$). Both designs are standardized by their training
moments before the kernel is formed, as the campaign's heads are, and the
check is at one seed. Read as diagnostics, they say that the learned features hand
the same kernel a target it interpolates with a much smaller fitted RKHS norm, and a somewhat smaller pointwise design factor at the median scale; they do not identify a causal
decomposition into a fixed design factor and a moving norm.

The associated feature-kernel function space has a standard pullback description.
Writing $\varphi$ for the feature map and $k$ for the base kernel, the deep
kernel head uses $k_\varphi(x,x')=k(\varphi(x),\varphi(x'))$, the construction
studied as deep kernel learning \citep{wilson2016deep,calandra2016manifold}.

\begin{proposition}[feature pullback]\label{prop:pullback}
The native space of $k_\varphi$ is $\mathcal H_{k_\varphi}=\{f\circ\varphi:
f\in\mathcal H_k\}$ with
$\|g\|_{\mathcal H_{k_\varphi}}=\min\{\|f\|_{\mathcal H_k}: f\circ\varphi=g
\text{ on }\mathcal U\}$. In particular, if the target factors through the
features as $G=h\circ\varphi$ with $h\in\mathcal H_k$, then
$\|G\|_{\mathcal H_{k_\varphi}}\le\|h\|_{\mathcal H_k}$, and the
optimal-recovery bound of Theorem~\ref{thm:or} for the feature-kernel
regression is governed by $\|h\|_{\mathcal H_k}$ rather than by the norm of
$G$ in the raw-input space.
\end{proposition}
\noindent\emph{Proof.} See Supplement~\ref{proof:prop:pullback}.


The statement is for scalar $f$; for the vector-valued targets of the
campaigns it is applied coordinatewise, with the norm the root sum of squares
over output coordinates, which is the norm Theorem~\ref{thm:or} uses.

The proof (Supplement~\ref{app:proofs}) is the standard pullback identity for
reproducing kernels \citep{saitoh1997integral,paulsen2016introduction}. The identity specifies the feature-kernel RKHS and the minimum extension norm. It provides no comparison with the raw-input RKHS without additional assumptions. In particular, a linear feature map is itself an input metric transformation. The measured fitted-norm reduction is evidence about these fitted kernels at the reported scales and nuggets, rather than a theorem about representation learning.

\subsection{An error bound for the residual kernel correction}
\label{sec:theory-or}

The final stage of the pipeline is kernel ridge regression of the ensemble
residual. The bound below is a vector-valued, residual form of the classical
power-function estimate for kernel interpolation
\citep{wendland2004scattered} and of the optimal-recovery viewpoint of
\citet{owhadi2019operator}; we include the short argument in
Supplement~\ref{app:proofs} to keep the two factors explicit. Fix the mean
$m:\mathcal U\to\mathcal V$ (in the pipeline, the stacked ensemble; in the
statement below $m$ is any realized map; no sample independence is needed for this deterministic statement) and
write $r=G-m$ for the residual operator with components $r_j$, $j=1,\dots,q$. Let $k$ be a positive definite kernel on $\mathcal U$ with
RKHS $\mathcal H_k$, and assume $r_j\in\mathcal H_k$ for all $j$ with
\[
\|r\|_K^2:=\sum_{j=1}^q\|r_j\|_{\mathcal H_k}^2<\infty .
\]
Given inputs $X=(u_1,\dots,u_n)$, Gram matrix $K=k(X,X)$ and nugget
$\lambda\ge 0$, the correction is
$\widehat r_\lambda(u)=k(u,X)\,(K+n\lambda I)^{-1}R$, $R_{nj}=r_j(u_n)$, and the
corrected predictor is $m+\widehat r_\lambda$. Let
\begin{equation}\label{eq:power}
P_\lambda(u)^2\;=\;k(u,u)-k(u,X)\,(K+n\lambda I)^{-1}k(X,u)
\end{equation}
denote the Gaussian process posterior variance with the same nugget (for
$\lambda=0$ this is the classical power function of the point set $X$).

\begin{theorem}[conditional power-function bound]\label{thm:or}
For every $u\in\mathcal U$ and every $\lambda\ge0$ (for $\lambda=0$ take $K$
positive definite, or read $K^{-1}$ as the pseudoinverse),
\[
\big\|G(u)-m(u)-\widehat r_\lambda(u)\big\|_2\;\le\;\|G-m\|_K\;
\widetilde P_\lambda(u)\;\le\;\|G-m\|_K\;P_\lambda(u),
\]
where
$\widetilde P_\lambda(u)^2=P_\lambda(u)^2-n\lambda\,\|(K+n\lambda
I)^{-1}k(X,u)\|_2^2$. The first inequality is sharp: for every $u$ there is
a residual operator of norm $\|G-m\|_K$ at which it holds with equality, so
$\|G-m\|_K\widetilde P_\lambda(u)$ is the exact worst-case error of the correction
over the ball $\{r:\|r\|_K\le\|G-m\|_K\}$.
\end{theorem}
\noindent\emph{Proof.} See Supplement~\ref{proof:thm:or}.


\paragraph{The recorded correction labels.}
Let $M_{\rm tr}\in\mathbb R^{n\times q}$ denote the training-row mean actually used by the implementation, and let $m_{\rm full}$ be the query-time mean. The kernel channel uses foldwise fits with pooled target centering; neural channels are in-sample, and the refiner's kernel input also changes at deployment. Thus $M_{\rm tr}$ need not be the restriction of a single cross-fitted function. Define
\[
 c_\lambda(u)=(K+n\lambda I)^{-1}k(X,u),\qquad
 \Delta=m_{\rm full}(X)-M_{\rm tr},\qquad
 R_{\rm tr}=G(X)-M_{\rm tr}.
\]
The implemented predictor is $m_{\rm full}(u)+c_\lambda(u)^\top R_{\rm tr}$. If $G-m_{\rm full}\in H_k^q$, Theorem~\ref{thm:or} and the triangle inequality give
\begin{align}
\|G(u)-m_{\rm full}(u)-c_\lambda(u)^\top R_{\rm tr}\|_2
&\le \|G-m_{\rm full}\|_K\widetilde P_\lambda(u)
       +\|c_\lambda(u)^\top\Delta\|_2,\label{eq:mismatch}\\
\|c_\lambda(u)^\top\Delta\|_2&\le\|c_\lambda(u)\|_2\|\Delta\|_F.\nonumber
\end{align}
Indeed, $R_{\rm tr}=(G-m_{\rm full})(X)+\Delta$; substituting this identity and applying the theorem proves the display. This proof is included here. The mismatch has not been measured, so the single-mean theorem is not an evaluated certificate for the reported correction.

The RKHS assumption is also unverified for these targets and the norm is unknown. For labels sampled from one compatible $r\in H_k^q$, ridge contraction gives $\|\widehat r_\lambda\|_K\le\|r\|_K$. With the recorded mixed training labels, the fitted norm need not lower-bound $\|G-m_{\rm full}\|_K$. Table~\ref{tab:norms} therefore describes fitted functions only: its squared-norm ratio is about 4900 before energy normalization and 8.8 after it. Neither ratio proves a corresponding reduction in an unknown residual norm or smoothness. The design factor is a worst-case error factor and need not rank realized errors well; Section~\ref{sec:uq} measures that ranking directly.

The design factor cannot be improved by a better estimator either. Among
all rules that see the same training residuals, the interpolant is optimal
in the worst case over the ball, and the price of a positive nugget in that
worst case is explicit.

\begin{theorem}[minimax optimality over the native-space ball]\label{thm:minimax}
Fix a design $X$ with $K$ positive definite, a point $u$ and a radius
$\rho>0$. For every map $\Psi:\mathbb R^{n\times q}\to\mathbb R^q$ that
estimates $r(u)$ from the training residuals $r(X)$,
\[
\sup_{\|r\|_K\le\rho}\big\|r(u)-\Psi\big(r(X)\big)\big\|_2\;\ge\;\rho\,P_0(u),
\]
and the kernel interpolant $\Psi(R)=k(u,X)K^{-1}R$ attains the bound. For
$\lambda>0$ the worst case of the ridge correction over the same ball is
$\rho\,\widetilde P_\lambda(u)$, and, with $k_u=k(X,u)$ and
$A_\lambda=K+n\lambda I$,
\[
P_0(u)^2\;\le\;\widetilde P_\lambda(u)^2\;\le\;P_\lambda(u)^2,\qquad
\widetilde P_\lambda(u)^2-P_0(u)^2=(n\lambda)^2\,k_u^\top
A_\lambda^{-2}K^{-1}k_u ,
\]
so the nugget's cost in the worst-case constant is a computable quantity that
vanishes as $\lambda\to0$.
\end{theorem}
\noindent\emph{Proof.} See Supplement~\ref{proof:thm:minimax}.


The lower bound is the classical optimal-recovery argument
\citep{owhadi2019operator}: two residuals of norm $\rho$ that vanish on the
design and take opposite values at $u$ are indistinguishable from the data,
so any rule errs by at least $\rho P_0(u)$ on one of them. What the theorem
adds for the pipeline is a reading of the reported quantity: at $\lambda=0$
the worst-case constant $\|G-m\|_K\,P_0(u)$ over the native-space ball is
the smallest any rule can attain from these data, and at the validated
nugget the gap to it is the displayed correction term, evaluable from the
Gram matrix alone. The statement is about the worst case over a ball whose
radius is not known; it is not an operational error bar for the shipped
surrogate, which is the split-conformal set of the next subsection.

\subsection{Distribution-free coverage for the reported band}
\label{sec:theory-conf}
Theorem~\ref{thm:or} controls the error by $\|G-m\|_K\,P_\lambda(u)$, but
the constant $\|G-m\|_K$ is not known a priori and
Section~\ref{sec:uq} shows that $P_\lambda$ alone ranks the error weakly.
What the surrogate reports as uncertainty is therefore a rescaling
$q\,P_\lambda(u)$, whose multiplier is the split-conformal quantile of the
scores $\|e(\tilde u_i)\|_2/P_\lambda(\tilde u_i)$ on a calibration split of
the test block, drawn after every selection in the pipeline has been made
on the validation split; Section~\ref{sec:uq} reports this band beside a
constant-width band calibrated the same way from the raw scores
$\|e(\tilde u_i)\|_2$. Coverage does not require the RKHS bound or a useful error ranking. For a predictor and positive scale fixed independently of exchangeable calibration and future inputs, Proposition~\ref{prop:conf} gives marginal coverage at least $1-\alpha$. Its proof is in Supplement~\ref{app:proofs}. The upper bound $1-\alpha+1/(m+1)$ additionally requires almost surely distinct scores. Define the conformal quantile using the $k=\lceil(1-\alpha)(m+1)\rceil$-th order statistic of the $m$ scores augmented by $+\infty$, which also covers $k=m+1$.

If, conditional on the fitted predictor and scale, the calibration and evaluation scores are i.i.d. with a continuous distribution and $k\le m$, the random conditional coverage over calibration draws has law $\mathrm{Beta}(k,m+1-k)$. The number covered among $N$ fresh evaluation cases is Beta--binomial. Remark~\ref{rem:covlaw} states and derives this law in Supplement~\ref{supp:conf}. At $m=1000$ and $N=19000$ its nominal-90\% central reference interval is $[88.0\%,91.8\%]$. A random split alone does not establish continuity or undo study-level adaptation to an already inspected test block. The recorded coverage experiment is compared with this reference law under those conditions; it is not a guarantee conditional on each input or on arbitrary prior test-driven selection.
